\section{Aprendizaje Supervisado}
\label{sec:supervisado}

\subsection*{Definición del problema predictivo}

\textbf{Variable objetivo:} \texttt{total\_ac\_power} (potencia AC total de la planta, suma de los 24 inversores), de tipo continuo $\Rightarrow$ problema de \textbf{regresión}.

\textbf{Variables de entrada (98):} corriente DC (24), voltaje DC (23), corriente AC (24) y voltaje AC (24) de cada inversor (95 señales eléctricas), más las variables temporales \texttt{hour}, \texttt{month} y \texttt{year}. Se excluyen del conjunto de entrada las columnas de potencia AC por inversor (son los componentes directos de la variable objetivo) y las variables agregadas de planta usadas solo en el EDA.

\textbf{Conjunto de datos:} 309\,308 registros con generación $>0$ (se excluyen las horas nocturnas, ver Sección \ref{sec:preprocesamiento}), particionados 70/15/15 en entrenamiento/validación/prueba.

\subsection*{Selección de modelos}

Se seleccionaron tres modelos de complejidad creciente para poder determinar, mediante evidencia comparativa y no por elección arbitraria, si la relación entre las señales de entrada y la potencia total requiere o no capacidad de modelado no lineal:

\begin{itemize}
    \item \textbf{Regresión Lineal} — baseline interpretable; asume relación lineal entre señales eléctricas y potencia total.
    \item \textbf{Random Forest Regressor} (150 árboles, profundidad máxima 14) — captura no linealidades e interacciones sin requerir escalamiento, robusto a valores atípicos.
    \item \textbf{Perceptrón Multicapa (MLP)} en \texttt{PyTorch} — red neuronal profunda (256$\to$128$\to$64$\to$1, ReLU, Dropout) entrenada con el optimizador \textbf{Adam} (combina momentum y escalamiento adaptativo tipo RMSProp) y pérdida \textbf{MSE}; capa de salida lineal (sin \texttt{Softmax}, reservado para clasificación) por tratarse de un problema de regresión continua.
\end{itemize}

\subsection*{Entrenamiento}

El MLP se entrenó hasta 60 épocas con \textit{early stopping} (paciencia de 10 épocas sin mejora en la pérdida de validación), \textit{gradient clipping} (norma máxima 1.0), un \textit{scheduler} \texttt{ReduceLROnPlateau} que reduce la tasa de aprendizaje cuando la validación se estanca, y restauración de los pesos correspondientes a la época de menor error de validación al finalizar el entrenamiento (en vez de conservar los pesos de la última época). El entrenamiento se detuvo por \textit{early stopping} en la época 16.

\begin{figure}[H]
\centering
\includegraphics[width=0.75\textwidth]{fig_curva_aprendizaje.png}
\caption{Curva de aprendizaje del MLP — MSE de entrenamiento y validación por época.}
\end{figure}

La curva evidencia una pérdida de validación inestable y con picos de gran magnitud en comparación con la de entrenamiento, un síntoma típico de sensibilidad del MLP a valores extremos aislados en los datos (Sección \ref{sec:no_supervisado} identifica picos de potencia anómalos en dos inversores). La restauración de los pesos de la mejor época mitiga, pero no elimina, el efecto de esta inestabilidad en el modelo final.

\subsection*{Evaluación y comparación de modelos}

\begin{table}[H]
\centering
\begin{tabular}{lrrr}
\toprule
\textbf{Modelo} & \textbf{RMSE (W)} & \textbf{MAE (W)} & \textbf{R\textsuperscript{2}} \\
\midrule
Regresión Lineal & 6.94  & 5.10  & 0.9992 \\
Random Forest    & 8.16  & 4.86  & 0.9988 \\
MLP (PyTorch)    & 51.33 & 44.52 & 0.9536 \\
\bottomrule
\end{tabular}
\caption{Comparación de modelos supervisados sobre el conjunto de prueba (15\% de los datos, no visto durante el entrenamiento)}
\end{table}

\begin{figure}[H]
\centering
\includegraphics[width=\textwidth]{fig_comparacion_modelos.png}
\caption{Potencia real vs. predicha en el conjunto de prueba, para los tres modelos.}
\end{figure}

\textbf{Modelo seleccionado:} \textbf{Regresión Lineal}, por presentar el menor RMSE (6.94 W) y el mayor R\textsuperscript{2} (0.9992) en el conjunto de prueba, superando tanto a Random Forest (RMSE 8.16 W) como al MLP (RMSE 51.33 W). Este resultado, a primera vista contraintuitivo (el modelo más simple superó a los dos más complejos), se explica en la Sección \ref{sec:discusion}: la potencia AC de cada inversor es aproximadamente el producto de su corriente y voltaje de salida, una relación que 98 variables lineales ya aproximan muy bien, dejando poco margen de mejora a modelos no lineales y exponiéndolos, en cambio, a mayor varianza de entrenamiento.

\subsection*{Interpretación de resultados — Importancia de variables}

\begin{figure}[H]
\centering
\includegraphics[width=0.85\textwidth]{fig_importancia_variables.png}
\caption{Top 15 variables por Permutation Importance sobre Random Forest.}
\end{figure}

La variable más importante por un margen amplio es \texttt{inv\_11\_dc\_current}, seguida de \texttt{inv\_20\_dc\_current}; el resto del top 15 está dominado por corrientes AC de distintos inversores. Ninguna variable de \textbf{voltaje} aparece entre las 15 más relevantes, confirmando lo observado en la Sección \ref{sec:relaciones}: el voltaje AC varía en un rango angosto y aporta poca información adicional una vez conocida la corriente, mientras que la corriente —directamente proporcional a la irradiancia solar instantánea— es el factor que realmente explica las variaciones de potencia.

\textbf{Significado práctico:} el modelo permite estimar, a partir de las lecturas instantáneas de corriente y voltaje de los inversores, la potencia total que la planta está entregando a la red — útil para validación cruzada de medidores, detección temprana de discrepancias entre lo esperado y lo medido, y como componente de un sistema de monitoreo en tiempo real. \textbf{Limitación:} el modelo predice la potencia en el mismo instante en que se miden las señales de entrada (no es un pronóstico a futuro); para eso se usa el análisis de forecasting de la Sección \ref{sec:forecasting}.
